\documentclass{article}
\usepackage{graphicx}
\usepackage{amsfonts}
\usepackage{amsmath}
\usepackage{amsthm}
\usepackage{tikz}

\theoremstyle{definition}
\newtheorem{theorem}{Theorem}[section]
\newtheorem{corollary}[theorem]{Corollary}
\newtheorem{proposition}[theorem]{Proposition}
\newtheorem{lemma}[theorem]{Lemma}
\newtheorem{definition}[theorem]{Definition}

\theoremstyle{remark}
\newtheorem{remark}[theorem]{Remark}
\newtheorem{example}[theorem]{Example}
\newtheorem{claim}[theorem]{Claim}

\title{Taylor and Maclaurin series}
\author{Sarah Liu}
\date{January 2026}

\begin{document}

\maketitle

\section{Taylor Polynomials}
\begin{definition}
    Let $a\in \mathbb{R}$. A power series centered at a is a function $f$ defined by an equation like 
    \[
    f(x)=\sum_{n=0}^{\infty}c_n(x-a)^n=c_0+c_1(x-a)+c_2(x-a)^2+\dots
    \]
    where $c_0,c_1,c_2,\dots\in\mathbb{R}$.
\end{definition}

\begin{theorem}
    Let $f(x)=\sum_{n=0}^{\infty}c_n(x-a)^n$ be a power series centered at $a\in\mathbb{R}$.
    \begin{enumerate}
        \item The domain of $f$ is an interval centered at $a$:$(a-R,a+R)$,$(a-R,a+R]$,$[a-R,a+R)$,$[a-R,a+R]$,$\mathbb{R}$, and $\{a\}$. We call this the \textbf{interval of convergence (IC)} of $f$, and $R$ the radius of convergence.
        \item In the interior of the IC, the series is absolutely convergent. In the exterior of the IC, the series is divergent. At the endpoints, anything may happen.
        \item In the interior of the IC, power series can be manipulate by basic operations.
    \end{enumerate}
\end{theorem}

\section{Taylor and Maclaurin series}
\[
\begin{aligned}
    f(x)&=\sum_{n=0}^{\infty}\frac{f^{(n)}(a)}{n!}(x-a)^n\\
    &=f(a)+\frac{f^{'}(a)}{1!}(x-a)+\frac{f^{''}(a)}{2!}(x-a)^2+\frac{f^{'''}(a)}{3!}(x-a)^3+\cdots
\end{aligned}
\]
\[
    f(x)=\sum_{n=0}^{\infty}\frac{f^{(n)}(0)}{n!}x^n=f(0)+\frac{f^{'}(0)}{1!}x+\frac{f^{''}(0)}{2!}x^2+\cdots
\]

\section{Four Examples}
\[
\begin{aligned}
e^x&=\sum_{n=0}^{\infty}\frac{x^n}{n!}=1+x+\frac{x^2}{2!}+\frac{x^3}{3!}+\cdots\\
\sin x&=\sum_{n=0}^{\infty}(-1)^n\frac{x^{2n+1}}{(2n+1)!}=x-\frac{x^3}{3!}+\frac{x^5}{5!}+\cdots\\
\cos x&=\sum_{n=0}^{\infty}(-1)^n\frac{x^{2n}}{(2n)!}=1-\frac{x^2}{2!}+\frac{x^4}{4!}-\frac{x^6}{6!}+\cdots\\
\frac{1}{1-x}&=\sum_{n=0}^{\infty}x^n=1+x+x^2+x^3\cdots
\end{aligned}
\]

\section{Remainder}

\begin{theorem}[Taylor's Inequality]
    If $|f^{(n+1)}(x)|\leq M$ for $|x-a|\leq d$, then the remainder $R_n(x)$ of the Taylor series satisfies the inequality
    \[
    |R_n(x)|\leq\frac{M}{(n+1)!}|x-a|^{n+1} \qquad \text{for }|x-a|\leq d
    \]
\end{theorem}

\begin{theorem}
    Let $f\in C^{n+1}$ is continuous on an interval containing $a$ and $x$, then
    \[
    f(x)=f(a)+\frac{f^{'}(a)}{1!}(x-a)+\frac{f^{''}(a)}{2!}(x-a)^2+\cdots + \frac{f^{(n)}(a)}{n!}(x-a)^n
    +\frac{1}{n!}\int_{t=a}^xf^{(n+1)}(t)(x-t)^n dt
    \]
\end{theorem}


\begin{theorem}[Lagrange's Remainder Theorem]
    Let $f\in C^{n+1}$ in an interval containing $a$ and $x$.
    Then
    \[
    f(x)-\left(\sum_{j=0}^n\frac{f^{(j)}(a)}{j!}(x-a)^j\right) = \frac{f^{(j)}(c)}{(n+1)!}(x-a)^{n+1}
    \]
    where $c$ is some number between $a$ and $x$.
\end{theorem}

\begin{proof}
    
\end{proof}

\end{document}